\documentclass[12pt,a4paper]{article}

% ------------------------------------------------
% Pacotes
% ------------------------------------------------
\usepackage[utf8]{inputenc}
\usepackage[T1]{fontenc}
\usepackage[brazil]{babel}
\usepackage{graphicx}
\usepackage[table]{xcolor}
\usepackage{array}
\usepackage{tabularx}
\usepackage{geometry}
\usepackage{fancyhdr} 
\usepackage{mathptmx} % Times New Roman 
\usepackage{enumitem}
\usepackage{amsfonts}
\usepackage{amsmath}
\usepackage{amssymb}
\usepackage{tikz}
\usetikzlibrary{positioning, shapes.geometric}

% TikZ Styles for ER Diagrams
\tikzset{
    block/.style={rectangle, draw, thick, minimum width=2cm, minimum height=1cm, fill=white},
    diamond_node/.style={diamond, draw, thick, minimum width=2cm, minimum height=1cm, fill=white, aspect=2},
    attribute/.style={ellipse, draw, thin, minimum width=1.5cm, minimum height=0.6cm, fill=white, font=\scriptsize},
    line/.style={draw, thick},
    specialization/.style={draw, circle, inner sep=1pt, font=\tiny, fill=white}
}

% ------------------------------------------------ 
% Configuração de página
% ------------------------------------------------
\geometry{
    a4paper,
    left=2cm,
    right=2cm,
    top=2cm,
    bottom=2cm,
    headheight=15pt
}

\pagestyle{fancy}
\fancyhf{}
\renewcommand{\headrulewidth}{0pt}
\setlength{\parindent}{0pt}

% Caixa para resposta
\newcommand{\resposta}[1]{
\noindent
\framebox[\textwidth]{
\begin{minipage}{0.98\textwidth}
\vspace{#1}
\end{minipage}
}
}

\newcolumntype{L}[1]{>{\raggedright\arraybackslash}p{#1}}

% ------------------------------------------------
% Documento
% ------------------------------------------------
\begin{document}

% ------------------------------------------------
% Cabeçalho
% ------------------------------------------------
\begin{tabularx}{\textwidth}{@{} L{4cm} X @{}}
    \raisebox{-0.5cm}{\includegraphics[width=4cm]{Universidade_de_Rio_Verde.jpg}} &
    \hfill
    \begin{tabular}{|>{\centering\arraybackslash}m{10cm}|>{\raggedright\arraybackslash}p{1.8cm}|}
        \hline
        \cellcolor{black}\color{white}\rule{0pt}{1.2cm}\textbf{LISTA DE EXERCÍCIOS: MODELAGEM DE DADOS} 
        & 
        \scriptsize Nota \tabularnewline
        \hline
    \end{tabular} 
\end{tabularx}
\vspace{0.3cm}

\begin{tabularx}{\textwidth}{|X|L{2cm}|}
\hline
\textbf{Curso} & \textbf{Campus} \\
\textit{Engenharia de Software} & \textit{Rio Verde} \\
\hline
\multicolumn{2}{|l|}{\textbf{Disciplina}} \\  
\multicolumn{2}{|l|}{ \textit{Banco de Dados I}} \\
\hline
\textbf{Nome do(a) acadêmico(a)} & \textbf{Valor} \\
& - \\
\hline
\textbf{Professor} & \textbf{Turma} \\
\textit{ Sergio Souza Novak} & \textit{A/B} \\
\hline
\textbf{Data} & \textbf{Código} \\
\textit{17/03/2026} & \textit{ESW} \\
\hline
\end{tabularx}

\vspace{0.5cm}

\textbf{Instruções:} Para cada um dos cenários abaixo, você deve:
\begin{enumerate}[label=\alph*)]
    \item Desenhar o \textbf{Diagrama Entidade-Relacionamento (Conceitual)}, contendo as 8 entidades, atributos principais e cardinalidades.
    \item Apresentar o \textbf{Esquema Relacional (Lógico)}, indicando Chaves Primárias (PK) e Chaves Estrangeiras (FK).
\end{enumerate}

\vspace{0.3cm}
\rule{\textwidth}{0.4pt}
\vspace{0.3cm}

% ------------------------------------------------
% EXERCÍCIO 1: E-COMMERCE
% ------------------------------------------------
\section*{Exercício 1: MarketplaceRV (E-commerce)}

Um novo marketplace precisa organizar suas vendas online. As regras de negócio são:
\begin{itemize}
    \item \textbf{Clientes e Pedidos:} O sistema registra Clientes (CPF, nome, email). Um cliente pode realizar vários Pedidos, mas cada pedido pertence a um único cliente.
    \item \textbf{Vendedores e Produtos:} Vendedores (CNPJ, nome\_loja) vendem diversos Produtos (ID, nome, preço). Um produto é vendido por apenas um vendedor.
    \item \textbf{Categorias:} Cada Produto obrigatoriamente pertence a uma única Categoria (ID, descrição), que pode agrupar vários produtos.
    \item \textbf{Itens de Pedido:} Um Pedido é composto por vários Produtos, e um Produto pode estar presente em vários pedidos diferentes. Deve-se registrar a quantidade e o preço unitário no momento da compra.
    \item \textbf{Pagamentos:} Cada Pedido gera exatamente um registro de Pagamento (ID, data, forma\_pgto, valor).
    \item \textbf{Cupons de Desconto:} O sistema possui Cupons (código, valor\_desconto). Um Cupom pode ser aplicado em vários Pedidos, mas cada pedido aceita no máximo um cupom.
\end{itemize}

\textbf{Entidades (8):} Cliente, Vendedor, Produto, Categoria, Pedido, Item\_Pedido, Pagamento, Cupom.

\resposta{10cm}

\newpage

% ------------------------------------------------
% EXERCÍCIO 2: BIBLIOTECA
% ------------------------------------------------
\section*{Exercício 2: BiblioRV (Gestão de Biblioteca)}

A biblioteca universitária deseja automatizar seu controle de acervo e empréstimos:
\begin{itemize}
    \item \textbf{Livros e Autores:} Livros (ISBN, título, ano) são escritos por um ou mais Autores (ID, nome, nacionalidade). Um autor pode ter escrito vários livros.
    \item \textbf{Editoras e Categorias:} Cada Livro é publicado por uma única Editora (ID, nome, cidade) e pertence a uma Categoria (ID, gênero) literária.
    \item \textbf{Exemplares:} Para cada Livro, a biblioteca possui vários Exemplares físicos (código\_de\_barras, estado\_conservação). Um exemplar refere-se a apenas um livro.
    \item \textbf{Leitores e Empréstimos:} Leitores (Matrícula, nome, tipo) podem realizar diversos Empréstimos de exemplares ao longo do semestre.
    \item \textbf{Registro de Empréstimo:} Cada Empréstimo registra a data de saída e a data prevista de devolução. Um empréstimo envolve exatamente um leitor e um único exemplar por vez.
    \item \textbf{Reservas:} Caso um Livro não tenha exemplares disponíveis, um Leitor pode fazer uma Reserva (ID, data\_reserva) para aquele livro. Um livro pode ter várias reservas na fila.
\end{itemize}

\textbf{Entidades (8):} Livro, Autor, Editora, Categoria, Exemplar, Leitor, Empréstimo, Reserva.

\resposta{12cm}

\newpage

% ------------------------------------------------
% EXERCÍCIO 3: ACADEMIA
% ------------------------------------------------
\section*{Exercício 3: FitRV (Gestão de Academia)}

Uma academia de ginástica precisa de um sistema para gerenciar treinos e mensalidades:
\begin{itemize}
    \item \textbf{Alunos e Planos:} Alunos (ID, nome, data\_nasc) matriculam-se em um único Plano (ID, nome\_plano, valor\_mensal) por vez. Um plano pode ter muitos alunos.
    \item \textbf{Instrutores e Modalidades:} Instrutores (CPF, nome, especialidade) ministram diversas Modalidades (ID, nome\_modalidade, sala). Cada modalidade tem um instrutor responsável.
    \item \textbf{Avaliações Físicas:} Um Aluno realiza várias Avaliações Físicas (ID, data, peso, \%\_gordura) periodicamente. Cada avaliação pertence a apenas um aluno.
    \item \textbf{Treinos e Exercícios:} Um Instrutor prescreve diversos Treinos (ID, objetivo, data\_inicio) para os alunos. Cada treino é composto por vários Exercícios (ID, nome\_exercicio, grupo\_muscular).
    \item \textbf{Montagem de Treino:} Um Exercício pode estar presente em vários treinos diferentes de alunos distintos. Deve-se registrar o número de séries e repetições para cada exercício no treino.
    \item \textbf{Equipamentos:} Exercícios utilizam Equipamentos (ID, nome\_equip, data\_manutenção). Um exercício pode exigir um equipamento, mas um equipamento pode ser usado em vários exercícios.
\end{itemize}

\textbf{Entidades (8):} Aluno, Plano, Instrutor, Modalidade, Avaliação, Treino, Exercício, Equipamento.

\resposta{12cm}

\end{document}
